% ─────────────────────────────────────────────────────────────────────────────
%  Chương 1 — Giới thiệu
%  Hướng dẫn của Khoa: 1--2 trang. Nêu công việc được giao, phát biểu bài toán,
%  khảo sát công nghệ kèm ưu nhược điểm, và lý do chọn công nghệ đã dùng.
% ─────────────────────────────────────────────────────────────────────────────
\chapter{Giới thiệu}
\label{ch:gioithieu}

\section{Công việc được giao}
\label{sec:congviec}

\todo{Chỉnh lại đoạn này theo đúng bối cảnh công ty bạn thực tập. Phần dưới viết
sẵn theo giả định phổ biến, giữ nguyên được thì giữ.}

Đơn vị thực tập giao cho tôi xây dựng một hệ thống giám sát camera có khả năng
tự động đếm lượng người ra vào một khu vực, thay cho cách đếm thủ công hoặc dùng
cảm biến hồng ngoại đặt tại cửa. Yêu cầu đặt ra gồm ba phần: xem được hình ảnh
trực tiếp từ nhiều camera cùng lúc, tự động nhận diện và đếm người qua một ranh
giới do người dùng tự vẽ, và tra cứu lại được các sự kiện đã xảy ra.

Cảm biến hồng ngoại tại cửa rẻ và dễ lắp, song đếm sai khi hai người đi sát nhau
và không phân biệt được chiều vào với chiều ra. Camera thì hầu hết cơ sở đã lắp
sẵn. Bài toán vì thế chuyển thành khai thác dữ liệu hình ảnh có sẵn thay vì đầu
tư thêm phần cứng.

\section{Phát biểu bài toán}
\label{sec:baitoan}

Cho một luồng video từ camera, hệ thống cần:

\begin{enumerate}[leftmargin=*, itemsep=2pt]
  \item Phát hiện vị trí của người trong từng khung hình.
  \item Gán cho mỗi người một mã định danh và duy trì mã đó khi người di chuyển
        giữa các khung hình liên tiếp.
  \item Đếm số lượt vượt qua một đoạn thẳng do người dùng vẽ, tách riêng hai
        chiều Vào và Ra.
  \item Lưu lại ảnh chụp và thời điểm của từng lượt để tra cứu về sau.
\end{enumerate}

Ràng buộc thực tế: hệ thống chạy trên một máy tính thông thường, không có máy
chủ GPU chuyên dụng; xử lý phải theo kịp luồng video thời gian thực; và người
dùng cuối là nhân viên vận hành, không phải kỹ sư, nên mọi cấu hình phải làm được
qua giao diện đồ họa.

\section{Khảo sát công nghệ}
\label{sec:khaosat}

\subsection{Mô hình phát hiện đối tượng}

Ba hướng tiếp cận phổ biến cho bài toán phát hiện đối tượng được cân nhắc, tóm
tắt trong Bảng~\ref{tab:sosanh-detector}.

\begin{table}[H]
\centering
\caption{So sánh các họ mô hình phát hiện đối tượng cho bài toán đếm người theo
thời gian thực.}
\label{tab:sosanh-detector}
\small
\begin{tabularx}{\textwidth}{@{}lXX@{}}
\toprule
\textbf{Mô hình} & \textbf{Ưu điểm} & \textbf{Nhược điểm} \\
\midrule
Faster R-CNN & Độ chính xác cao, đặc biệt với vật thể nhỏ &
Hai giai đoạn nên chậm, khó đạt thời gian thực trên CPU \\
\addlinespace
DETR và biến thể & Bỏ được bước NMS, kiến trúc gọn về mặt khái niệm &
Cần nhiều dữ liệu và thời gian huấn luyện, hệ sinh thái triển khai còn mỏng \\
\addlinespace
YOLO (v8) & Một giai đoạn, tốc độ cao, trọng số huấn luyện sẵn trên COCO,
thư viện Ultralytics hỗ trợ đầy đủ &
Độ chính xác thấp hơn Faster R-CNN với vật thể rất nhỏ hoặc bị che khuất nhiều \\
\bottomrule
\end{tabularx}
\end{table}

YOLOv8 được chọn. Lý do quyết định là bộ trọng số huấn luyện sẵn trên COCO đã
chứa lớp \code{person} cùng các lớp phương tiện, nên hệ thống chạy được ngay mà
không cần gán nhãn và huấn luyện lại. Với một kỳ thực tập kéo dài vài tuần, việc
xây dựng tập dữ liệu riêng sẽ chiếm gần hết quỹ thời gian và không còn chỗ cho
phần ứng dụng.

\subsection{Thuật toán bám vết}

Phát hiện đơn thuần thì chưa đếm được. Muốn biết một người đã đi qua vạch hay
chưa, cần theo dõi người đó qua nhiều khung hình liên tiếp.

SORT chạy rất nhanh nhưng đổi mã định danh liên tục mỗi khi đối tượng bị che
khuất. DeepSORT khắc phục bằng cách so khớp thêm đặc trưng ngoại hình, đổi lại
phải chạy thêm một mạng trích đặc trưng cho từng hộp, chi phí tăng đáng kể.
ByteTrack chọn hướng khác: tận dụng cả những hộp có điểm tin cậy thấp mà SORT
loại bỏ, nhờ đó giữ được vết qua các đoạn che khuất ngắn mà không cần mạng phụ
nào. Thư viện \en{Supervision} đã đóng gói sẵn ByteTrack cùng các tiện ích vẽ chú
thích, nên chi phí tích hợp gần như bằng không.

\subsection{Truyền hình ảnh về trình duyệt}

WebRTC cho độ trễ thấp nhất nhưng đòi hỏi máy chủ báo hiệu, xử lý STUN/TURN và
mã hóa lại luồng. HLS ổn định qua Internet, đổi lại độ trễ vài giây do phải chia
đoạn. MJPEG thì chỉ là chuỗi ảnh JPEG nối tiếp trong một phản hồi HTTP, hiển thị
được bằng thẻ \code{<img>} thông thường.

Hệ thống chạy trong mạng nội bộ với số camera nhỏ, nên MJPEG là lựa chọn hợp lý:
độ trễ dưới nửa giây, không cần thư viện phía trình duyệt, và quan trọng hơn cả
là ảnh gửi đi đã được vẽ sẵn hộp giới hạn cùng vạch đếm ở phía máy chủ. Mọi máy
trạm nhìn thấy đúng một kết quả. Cái giá phải trả là băng thông cao hơn, điều
này chấp nhận được trong mạng LAN.

\subsection{Cơ sở dữ liệu}

PostgreSQL mạnh khi nhiều tiến trình cùng ghi và dữ liệu lớn, nhưng kéo theo một
dịch vụ phải cài đặt và vận hành riêng. SQLite chạy ngay trong tiến trình ứng
dụng, không cần cấu hình, và ở chế độ WAL cho phép nhiều luồng đọc song song với
một luồng ghi. Khối lượng ghi của hệ thống chỉ vài sự kiện mỗi phút cho mỗi
camera, thấp hơn nhiều so với ngưỡng SQLite bắt đầu đuối. Chọn SQLite giúp toàn
bộ hệ thống triển khai bằng một lệnh duy nhất.

\section{Công nghệ sử dụng}
\label{sec:congnghe}

\begin{table}[H]
\centering
\caption{Thành phần công nghệ của hệ thống và vai trò của từng thành phần.}
\label{tab:congnghe}
\small
\begin{tabularx}{\textwidth}{@{}llX@{}}
\toprule
\textbf{Lớp} & \textbf{Công nghệ} & \textbf{Vai trò} \\
\midrule
Xử lý ảnh    & OpenCV            & Đọc luồng video, vẽ chú thích, mã hóa JPEG và ghi tệp MP4 \\
Phát hiện    & YOLOv8 (Ultralytics) & Xác định vị trí người và phương tiện trong khung hình \\
Bám vết      & ByteTrack (Supervision) & Gán và duy trì mã định danh cho từng đối tượng \\
Máy chủ      & FastAPI, Uvicorn  & Cung cấp REST API và luồng MJPEG \\
Lưu trữ      & SQLite (chế độ WAL) & Ghi sự kiện, số đếm, chỉ mục đoạn video \\
Giao diện    & React, TypeScript, Vite & Ứng dụng một trang cho người vận hành \\
Biểu đồ      & Recharts          & Vẽ biểu đồ lưu lượng theo giờ \\
\bottomrule
\end{tabularx}
\end{table}

Python được chọn cho phần xử lý vì toàn bộ hệ sinh thái thị giác máy tính đều
xoay quanh ngôn ngữ này. TypeScript cho phần giao diện giúp bắt lỗi sai kiểu dữ
liệu ngay khi biên dịch, thay vì đợi đến lúc chạy mới phát hiện API trả về thiếu
trường.
